\subsection{Adjunction lemma}%\label{s1.4}

The following statement is standard when dealing with categories with enough projective or injective objects.
We failed to f\/ind it in the literature in the generality we need, so we provide a~proof without the use of projective or
injective objects.

\begin{proposition}[adjunction lemma]
\label{adjlem}
Let $\mathcal{A}$ and $\mathcal{B}$ be two Abelian categories and $(\mathrm{F},\mathrm{G})$ an adjoint pair of exact
functors $\mathrm{F}:\mathcal{A}\to \mathcal{B}$ and $\mathrm{G}:\mathcal{B}\to \mathcal{A}$.
Then for every $d\in\mathbb{Z}_{\geq 0}$, $N\in \mathcal{A}$ and $M\in \mathcal{B}$ there are isomorphisms
\begin{gather*}
\mathrm{Ext}^d_{\mathcal{B}}(\mathrm{F}(N),M)\cong \mathrm{Ext}^d_{\mathcal{A}}(N,\mathrm{G}(M))
\end{gather*}
natural in both~$N$ and~$M$.
\end{proposition}

\begin{proof}
Applying $\mathrm{F}$ to an exact sequence
\begin{gather*}
\xymatrix@C=7mm{
0\ar[r]&\mathrm{G}(M)\ar[r]&X_1\ar[r]&X_2\ar[r]&\cdots\ar[r]&X_{d-1}\ar[r]&X_d\ar[r]&N\ar[r]&0
}
\end{gather*}
in $\mathcal{A}$, gives an exact sequence
\begin{gather*}
\xymatrix@C=5mm{
0\ar[r]&\mathrm{F}\mathrm{G}(M)\ar[r]&\mathrm{F}(X_1)\ar[r]&
\mathrm{F}(X_2)\ar[r]&\cdots\ar[r]&\mathrm{F}(X_{d-1})\ar[r]&\mathrm{F}(X_d)\ar[r]&\mathrm{F}(N)\ar[r]&0
}
\end{gather*}
in $\mathcal{B}$.
Denote by $\mathrm{K}$ the kernel of the adjunction natural transformation $\mathrm{F}\mathrm{G}\to
\mathrm{Id}_{\mathcal{B}}$.
Then we have the following commutative diagram
\begin{gather}
\label{eq11}
\begin{split}
& \xymatrix@C=5mm{
0\ar[r]&\mathrm{K}(M)\ar[r]\ar@{^{(}->}[d]&\mathrm{K}(M)\ar[r]\ar@{^{(}->}[d]&
0\ar[d]\ar[r]&\cdots\ar[r]&0\ar[d]\ar[r]&0\ar[d]\ar[r]&0\ar[d]\ar[r]&0\\
0\ar[r]&\mathrm{F}\mathrm{G}(M)\ar[r]\ar@{->>}[d]&\mathrm{F}(X_1)\ar[r]\ar@{->>}[d]&
\mathrm{F}(X_2)\ar[r]\ar@{=}[d]&\cdots\ar[r]&\mathrm{F}(X_{d-1})\ar[r]\ar@{=}[d]&
\mathrm{F}(X_d)\ar[r]\ar@{=}[d]&\mathrm{F}(N)\ar[r]\ar@{=}[d]&0\\
0\ar[r]&M'\ar[r]\ar@{^{(}->}[d]&X'\ar[r]\ar@{^{(}->}[d]&
\mathrm{F}(X_2)\ar[r]\ar@{=}[d]&\cdots\ar[r]&\mathrm{F}(X_{d-1})\ar[r]\ar@{=}[d]&
\mathrm{F}(X_d)\ar[r]\ar@{=}[d]&\mathrm{F}(N)\ar[r]\ar@{=}[d]&0\\
0\ar[r]&M\ar[r]&X''\ar[r]&
\mathrm{F}(X_2)\ar[r]&\cdots\ar[r]&\mathrm{F}(X_{d-1})\ar[r]&\mathrm{F}(X_d)\ar[r]&\mathrm{F}(N)\ar[r]&0
}
\end{split}
\end{gather}
with exact rows.
Here the homomorphism from the f\/irst to the second row is given by the natural embedding $\mathrm{K}\hookrightarrow
\mathrm{F}\mathrm{G}$ and the third row is just the corresponding cokernel with the morphism from the second to the
third row being the canonical projection.
In particular, $M':=\mathrm{F}\mathrm{G}(M)/\mathrm{K}(M)$.
The homomorphism from $M'$ to~$M$ is the natural inclusion (coming from the def\/inition of $\mathrm{K}$) and, f\/inally,
$X''$ is def\/ined as the push-out and the map to $\mathrm{F}(X_2)$ is given by the universal property of push-outs.
Functoriality of the construction yields a~group homomorphism
\begin{gather*}
\Phi: \ \mathrm{Ext}^d_{\mathcal{A}}(N,\mathrm{G}(M))\to \mathrm{Ext}^d_{\mathcal{B}}(\mathrm{F}(N),M).
\end{gather*}

Applying $\mathrm{G}$ to an exact sequence
\begin{gather*}
\xymatrix@C=7mm{
0\ar[r]&M\ar[r]&Y_1\ar[r]&Y_2\ar[r]&\cdots\ar[r]&Y_{d-1}\ar[r]&Y_d\ar[r]&\mathrm{F}(N)\ar[r]&0
}
\end{gather*}
in $\mathcal{B}$, gives an exact sequence
\begin{gather*}
\xymatrix@C=5mm{
0\ar[r]&\mathrm{G}(M)\ar[r]&\mathrm{G}(Y_1)\ar[r]&\mathrm{G}(Y_2)\ar[r]&
\cdots\ar[r]&\mathrm{G}(Y_{d-1})\ar[r]&\mathrm{G}(Y_d)\ar[r]&\mathrm{G}\mathrm{F}(N)\ar[r]&0
}
\end{gather*}
in $\mathcal{A}$.
Denote by $\mathrm{C}$ the cokernel of the adjunction natural transformation $\mathrm{Id}_{\mathcal{A}}\to
\mathrm{G}\mathrm{F}$.
Then we have the following commutative diagram
\begin{gather*}%\label{eq12}
\xymatrix@C=5mm{
0\ar[r]&\mathrm{G}(M)\ar[r]\ar@{=}[d]&\mathrm{G}(Y_1)\ar[r]\ar@{=}[d]&\mathrm{G}(Y_2)\ar[r]\ar@{=}[d]&
\cdots\ar[r]&\mathrm{G}(Y_{d-1})\ar[r]\ar@{=}[d]&
Y''\ar[r]\ar@{->>}[d]&N\ar[r]\ar@{->>}[d]&0\\
0\ar[r]&\mathrm{G}(M)\ar[r]\ar@{=}[d]&\mathrm{G}(Y_1)\ar[r]\ar@{=}[d]&\mathrm{G}(Y_2)\ar[r]\ar@{=}[d]&
\cdots\ar[r]&\mathrm{G}(Y_{d-1})\ar[r]\ar@{=}[d]&
Y'\ar[r]\ar@{^{(}->}[d]&N'\ar[r]\ar@{^{(}->}[d]&0\\
0\ar[r]&\mathrm{G}(M)\ar[r]\ar[d]&\mathrm{G}(Y_1)\ar[r]\ar[d]&\mathrm{G}(Y_2)\ar[r]\ar[d]&
\cdots\ar[r]&\mathrm{G}(Y_{d-1})\ar[r]\ar[d]&
\mathrm{G}(Y_d)\ar[r]\ar@{->>}[d]&\mathrm{G}\mathrm{F}(N)\ar[r]\ar@{->>}[d]&0\\
0\ar[r]&0\ar[r]&0\ar[r]&0\ar[r]&
\cdots\ar[r]&0\ar[r]&\mathrm{C}(N)\ar[r]&\mathrm{C}(N)\ar[r]&0
}
\end{gather*}
with exact rows.
Here the homomorphism from the third to the last row is given by the natural projection of
$\mathrm{G}\mathrm{F}\twoheadrightarrow \mathrm{C}$ and the second row is just the corresponding kernel with the
morphism from the second to the third row being the canonical injection.
In particular, $\mathrm{C}(N):=\mathrm{G}\mathrm{F}(M)/N'$.
The homomorphism from~$N$ to $N'$ is the natural surjection and, f\/inally, $Y''$ is def\/ined as the pullback and the map
from $\mathrm{G}(Y_{d-1})$ is given by the universal property of pullbacks.
Functoriality of the construction yields a~group homomorphism
\begin{gather*}
\Psi: \ \mathrm{Ext}^d_{\mathcal{B}}(\mathrm{F}(N),M)\to \mathrm{Ext}^d_{\mathcal{A}}(N,\mathrm{G}(M)).
\end{gather*}

Finally, we claim that~$\Phi$ and~$\Psi$ are inverses of each other.
Consider the following commutative diagram:
\begin{gather}
\label{eq14}
\begin{split}
& \xymatrix@C=3mm{
0\ar[r]&\mathrm{G}(M)\ar[r]\ar[d]&X_1\ar[r]\ar[d]&X_2\ar[r]\ar[d]&\cdots\ar[r]
&X_{d-1}\ar[r]\ar[d]&X_d\ar[r]\ar[d]&N\ar[r]\ar[d]&0\\
0\ar[r]&\mathrm{G}\mathrm{F}\mathrm{G}(M)\ar[r]\ar@{->>}[d]&\mathrm{G}\mathrm{F}(X_1)\ar[r]\ar@{->>}[d]&
\mathrm{G}\mathrm{F}(X_2)\ar[r]\ar@{=}[d]&\cdots\ar[r]&\mathrm{G}\mathrm{F}(X_{d-1})\ar[r]\ar@{=}[d]&
\mathrm{G}\mathrm{F}(X_d)\ar[r]\ar@{=}[d]&\mathrm{G}\mathrm{F}(N)\ar[r]\ar@{=}[d]&0\\
0\ar[r]&\mathrm{G}(M)\ar[r]&\mathrm{G}(X'')\ar[r]&
\mathrm{G}\mathrm{F}(X_2)\ar[r]&\cdots\ar[r]&\mathrm{G}\mathrm{F}(X_{d-1})\ar[r]&
\mathrm{G}\mathrm{F}(X_d)\ar[r]&\mathrm{G}\mathrm{F}(N)\ar[r]&0\\
0\ar[r]&\mathrm{G}(M)\ar[r]\ar@{=}[u]&\mathrm{G}(X'')\ar[r]\ar@{=}[u]&
\mathrm{G}\mathrm{F}(X_2)\ar[r]\ar@{=}[u]&\cdots\ar[r]&\mathrm{G}\mathrm{F}(X_{d-1})\ar[r]\ar@{=}[u]&
Y''\ar[r]\ar[u]&N\ar[r]\ar[u]&0
}
\end{split}
\end{gather}
Here the second row is obtained from the f\/irst one by applying $\mathrm{G}\mathrm{F}$ and the homomorphism from the
f\/irst to the second row is given by adjunction $\mathrm{Id}_{\mathcal{A}}\to\mathrm{G}\mathrm{F}$.
The second and the third rows and the homomorphism between them are given by applying the exact functor $\mathrm{G}$ to
the two middle rows of~\eqref{eq11}.
Note that, by construction, application of $\mathrm{G}$ identif\/ies the two last rows of~\eqref{eq11}, that is
$\mathrm{G}(M')\cong \mathrm{G}(M)$ and $\mathrm{G}(X')\cong \mathrm{G}(X'')$ (this follows from surjectivity of the
natural transformation $\mathrm{G}\mathrm{F}\mathrm{G}\to \mathrm{G}$ given by adjunction).
Consequently, from the adjunction identities we have that the composition of the maps in the f\/irst column is the
identity on $\mathrm{G}(M)$.
Finally, the last row and the homomorphism to the third row are given by the def\/inition of~$\Psi$.
In particular, from the construction we have that the image of~$N$ in the second row (coming from the f\/irst row) and in
the third row (coming from the last row) coincide.
Hence diagram~\eqref{eq14} gives rise to a~diagram
\begin{gather*}
\xymatrix@C=3mm{
0\ar[r]&\mathrm{G}(M)\ar[r]\ar@{=}[d]&X_1\ar[r]\ar[d]&X_2\ar[r]\ar[d]&\cdots\ar[r]
&X_{d-1}\ar[r]\ar[d]&X_d\ar[r]\ar[d]&N\ar[r]\ar[d]&0\\
0\ar[r]&\mathrm{G}(M)\ar[r]&\mathrm{G}(X'')\ar[r]&
\mathrm{G}\mathrm{F}(X_2)\ar[r]&\cdots\ar[r]&\mathrm{G}\mathrm{F}(X_{d-1})\ar[r]&
Q'\ar[r]&N'\ar[r]&0\\
0\ar[r]&\mathrm{G}(M)\ar[r]\ar@{=}[u]&\mathrm{G}(X'')\ar[r]\ar@{=}[u]&
\mathrm{G}\mathrm{F}(X_2)\ar[r]\ar@{=}[u]&\cdots\ar[r]&\mathrm{G}\mathrm{F}(X_{d-1})\ar[r]\ar@{=}[u]&
Y''\ar[r]\ar[u]&N\ar[r]\ar[u]&0
}
\end{gather*}
with exact rows, where $N'$ is the image of~$N$ in $\mathrm{G}\mathrm{F}(N)$ and $Q'$ is the full preimage of $N'$.
Pulling back along the epimorphism $N\twoheadrightarrow N'$ gives a~commutative diagram
\begin{gather*}
\xymatrix@C=3mm{
0\ar[r]&\mathrm{G}(M)\ar[r]\ar@{=}[d]&X_1\ar[r]\ar[d]&X_2\ar[r]\ar[d]&\cdots\ar[r]
&X_{d-1}\ar[r]\ar[d]&X_d\ar[r]\ar[d]&N\ar[r]\ar@{=}[d]&0\\
0\ar[r]&\mathrm{G}(M)\ar[r]&\mathrm{G}(X'')\ar[r]&
\mathrm{G}\mathrm{F}(X_2)\ar[r]&\cdots\ar[r]&\mathrm{G}\mathrm{F}(X_{d-1})\ar[r]&
Q\ar[r]&N\ar[r]&0\\
0\ar[r]&\mathrm{G}(M)\ar[r]\ar@{=}[u]&\mathrm{G}(X'')\ar[r]\ar@{=}[u]&
\mathrm{G}\mathrm{F}(X_2)\ar[r]\ar@{=}[u]&\cdots\ar[r]&\mathrm{G}\mathrm{F}(X_{d-1})\ar[r]\ar@{=}[u]&
Y''\ar[r]\ar[u]&N\ar[r]\ar@{=}[u]&0
}
\end{gather*}
with exact rows.
The last diagram shows that the extensions given by the f\/irst and the last rows coincide, which proves that $\Psi\Phi$
is the identity map.

The claim that $\Phi\Psi$ is the identity map is proved similarly.
This completes the proof.
\end{proof}

\section{Gelfand--Zeitlin modules}\label{s2}

\textbf{Notation.} For a~Lie algebra $\mathfrak{a}$ we denote by $U(\mathfrak{a})$ its universal enveloping algebra and
by $Z(\mathfrak{a})$ the center of $U(\mathfrak{a})$.

\subsection[Gelfand--Zeitlin subalgebra of $\mathfrak{gl}_n$]{Gelfand--Zeitlin subalgebra of $\boldsymbol{\mathfrak{gl}_n}$}%\label{s2.1}

For $k\in\mathbb{Z}_{>0}$ denote by $\mathfrak{g}_k$ the Lie algebra $\mathfrak{gl}_k(\mathbb{C})$.
Set $U_k=U(\mathfrak{g}_k)$ and $Z_k=Z(\mathfrak{g}_k)$.
We consider the usual chain
\begin{gather*}
\mathfrak{g}_1\subset \mathfrak{g}_2\subset \mathfrak{g}_3\subset \dots \subset \mathfrak{g}_n
\end{gather*}
of ``left upper corner'' embeddings of Lie algebras as depicted on the following picture:

\begin{center}
\begin{picture}(143.00,143.00)
\drawline(20,20)(140,20)
\drawline(20,20)(20,140)
\drawline(140,140)(20,140)
\drawline(140,140)(140,20)
\drawline(20,120)(40,120)
\drawline(40,140)(40,120)
\drawline(20,100)(60,100)
\drawline(60,140)(60,100)
\drawline(20,80)(80,80)
\drawline(80,140)(80,80)
\drawline(20,40)(120,40)
\drawline(120,140)(120,40)
\put(25,130){$\mathfrak{g}_1$}
\put(45,110){$\mathfrak{g}_2$}
\put(65,90){$\mathfrak{g}_3$}
\put(85,70){$\cdot$}
\put(87.5,67.5){$\cdot$}
\put(90,65){$\cdot$}
\put(100,50){$\mathfrak{g}_{n\text{-}1}$}
\put(125,30){$\mathfrak{g}_n$}
\end{picture}
\end{center}

\vspace{-6mm}

This gives rise to the chain
\begin{gather*}
U_1\subset U_2\subset U_3\subset \dots \subset U_n
\end{gather*}
of embeddings of associative algebras.
The subalgebra~$\Gamma$ of $U:=U_n$ generated by all centers~$Z_k$, where $k=1,2,\dots,n$, is called the {\em
Gelfand--Zeitlin} subalgebra.
We set $\mathfrak{g}:=\mathfrak{g}_n$.

The algebra~$\Gamma$ is obviously commutative, moreover,~$\Gamma$ is a~polynomial algebra in $\frac{n(n+1)}{2}$
variables (these can be taken to be generators of $Z_k$ for $k=1,2,\dots,n$), see~\cite{GZ} or~\cite[Chapter~X]{Zh}.
Considered as a~subalgebra of~$U$,~$\Gamma$ is a~{\em Harish-Chandra subalgebra} in the sense of~\cite[Section~1]{DFO},
which means that every f\/initely generated~$\Gamma$-subbimodule of~$U$ is already f\/initely generated both as a~left and
as a~right~$\Gamma$-module, see~\cite[Theorem~24]{DFO}.
Furthermore,~$U$ is free both as a~left and as a~right~$\Gamma$-module, see~\cite{Ov}.
Consequently, the usual induction and coinduction functors
\begin{gather*}
\mathrm{Ind}_{\Gamma}^U:=U\bigotimes_{\Gamma}{}_-
\qquad
\text{and}
\qquad
\mathrm{Coind}_{\Gamma}^U:=\mathrm{Hom}_{\Gamma}(U,{}_-)
\end{gather*}
are exact.

\subsection{Gelfand--Zeitlin modules}%\label{s2.2}

The category $\mathscr{G}\mathscr{Z}$ of {\em Gelfand--Zeitlin}-modules for~$U$ is def\/ined as the full subcategory of
$U\dash{mod}$ (the category of f\/initely generated~$U$-modules), which consists of those $M\in U\dash{mod}$ on which the
action of~$\Gamma$ is locally f\/inite, that is $\dim(\Gamma v)<\infty$ for all $v\in M$.
The category $\mathscr{G}\mathscr{Z}$ is a~Serre subcategory of $U\dash{mod}$ as~$U$ is Noetherian.

Write $\mathrm{Specm}(\Gamma)$ for the set of maximal ideals in~$\Gamma$.
As~$\Gamma$ is a~polynomial algebra, we have the decomposition
\begin{gather*}
\Gamma\dash{fmod}=\bigoplus_{\mathbf{m}\in\mathrm{Specm}(\Gamma)} \Gamma\dash{fmod}_{\mathbf{m}},
\end{gather*}
where $\Gamma\dash{fmod}_{\mathbf{m}}$ denotes the full subcategory of $\Gamma\dash{fmod}$ consisting of all objects
annihilated by some power of $\mathbf{m}$.
From this decomposition, consider the functor
\begin{gather*}
\mathrm{Ind}_{\Gamma,\mathbf{m}}^U: \ \Gamma\dash{fmod}_{\mathbf{m}}\to \mathscr{G}\mathscr{Z}
\end{gather*}
def\/ined as the restriction of $\mathrm{Ind}_{\Gamma}^U$ to $\Gamma\dash{fmod}_{\mathbf{m}}$.

Note that modules in $\mathscr{G}\mathscr{Z}$ are usually inf\/inite-dimensional and hence the usual restriction
$\mathrm{Res}_{\Gamma}^U$ ends up in $\Gamma\dash{lfmod}$ and not in $\Gamma\dash{fmod}$.
However, from~\cite[Corollary~5.3(a)]{FuOv2} it follows that for any $M\in \mathscr{G}\mathscr{Z}$ and
$\mathbf{m}\in\mathrm{Specm}(\Gamma)$ the space
\begin{gather*}
M_{\mathbf{m}}:=\big\{v\in \mathrm{Res}_{\Gamma}^U M: \mathbf{m}^iv=0~\text{for some}~i\big\}
\end{gather*}
is f\/inite-dimensional.
This allows us to def\/ine the functor
\begin{gather*}
\mathrm{Res}_{\Gamma,\mathbf{m}}^U: \ \mathscr{G}\mathscr{Z}\to \Gamma\dash{fmod}_{\mathbf{m}},
\end{gather*}
which sends~$M$ to $M_{\mathbf{m}}$ and is def\/ined on morphisms by restriction.
Then the usual adjunction between induction and restriction implies that $\mathrm{Ind}_{\Gamma,\mathbf{m}}^U$ is left
adjoint to $\mathrm{Res}_{\Gamma,\mathbf{m}}^U$.
The functor $\mathrm{Res}_{\Gamma,\mathbf{m}}^U$ is, in turn, left adjoint to the restriction
\begin{gather*}
\mathrm{Coind}_{\Gamma,\mathbf{m}}^U: \ \Gamma\dash{fmod}_{\mathbf{m}}\to \mathscr{G}\mathscr{Z}
\end{gather*}
of the coinduction functor to $\Gamma\dash{fmod}_{\mathbf{m}}$.

\subsection{The main result}%\label{s2.3}

Our main result in this section is the following:

\begin{theorem}
\label{thm3}
The category $\mathscr{G}\mathscr{Z}$ is extension full in $U\dash{Mod}$.
\end{theorem}

\begin{proof}
To prove Theorem~\ref{thm3}, we would like to apply Proposition~\ref{prop1} for $\mathcal{A}=U\dash{Mod}$,
$\mathcal{B}=\mathscr{G}\mathscr{Z}$ and~$\mathcal{B}_0$ being the full subcategory of $\mathcal{B}$ consisting of
all~$U$-modules isomorphic to $\mathrm{Ind}_{\Gamma}^U N$ for a~f\/inite-dimensional~$\Gamma$-module~$N$.
Taking the assumptions of Proposition~\ref{prop1} into account, Theorem~\ref{thm3} reduces to the following lemma:

\begin{lemma}
%\label{lem4}
Let $Q\in\mathscr{G}\mathscr{Z}$,~$N$ be a~finite-dimensional~$\Gamma$-module and $M=\mathrm{Ind}_{\Gamma}^U N$.
Then $\iota^d_{M,Q}$ is an isomorphism for every $d\in\mathbb{Z}_{\geq 0}$.
\end{lemma}

\begin{proof}
By additivity, we may assume $N\in \Gamma\dash{fmod}_{\mathbf{m}}$ for some $\mathbf{m}\in\mathrm{Specm}(\Gamma)$.
The image of the functor $\mathrm{Ind}_{\Gamma,\mathbf{m}}^U:\Gamma\dash{fmod}_{\mathbf{m}}\to \mathcal{A}$ belongs to
$\mathcal{B}$.
We show that for every $d\in\mathbb{Z}_{\geq 0}$, any $N\in \Gamma\dash{fmod}_{\mathbf{m}}$ and any
$Q\in\mathscr{G}\mathscr{Z}$ we have the isomorphisms
\begin{gather}
\label{eq1}
\mathrm{Ext}^d_{\mathcal{A}}\big(\mathrm{Ind}_{\Gamma,\mathbf{m}}^U N, Q\big)  \cong
\mathrm{Ext}^d_{\Gamma\dash{Mod}}\big(N,\mathrm{Res}_{\Gamma}^U Q\big),
\\
\label{eq2}
\mathrm{Ext}^d_{\Gamma\dash{Mod}}\big(N,\mathrm{Res}_{\Gamma}^U Q\big)  \cong
\mathrm{Ext}^d_{\Gamma\dash{mod}_{\mathbf{m}}}\big(N,\mathrm{Res}_{\Gamma,\mathbf{m}}^U Q\big),
\\
\label{eq3}
\mathrm{Ext}^d_{\Gamma\dash{mod}_{\mathbf{m}}}\big(N,\mathrm{Res}_{\Gamma,\mathbf{m}}^U Q\big)  \cong
\mathrm{Ext}^d_{\mathcal{B}}\big(\mathrm{Ind}_{\Gamma,\mathbf{m}}^U N, Q\big).
\end{gather}
Since, by construction, these three isomorphisms together with the morphism $\iota^d_{M,Q}$ will yield a~commutative
square, we get that $\iota^d_{M,Q}$ is an isomorphism.

Isomorphism~\eqref{eq2} follows from~\cite[Lemma~17]{CoMa}.
Isomorphisms~\eqref{eq1} and~\eqref{eq3} follow from adjunction lemma (Proposition~\ref{adjlem}).
Note that adjunction lemma applies here since~$U$ is free over~$\Gamma$.
\end{proof}

Theorem~\ref{thm3} now follows.
\end{proof}
